\documentclass[11pt]{article}
\usepackage{fontspec}
\setmainfont{Times New Roman}
\setsansfont{Arial}
\setmonofont{Consolas}
\usepackage{amsmath,amssymb}
\usepackage{geometry}
\usepackage{microtype}
\geometry{a4paper,margin=3cm}
\setlength{\parindent}{1.25em}
\setlength{\parskip}{0pt}
\emergencystretch=10em
\allowdisplaybreaks
\providecommand{\Kfield}{\mathfrak{K}}
\providecommand{\Ssys}{\mathfrak{S}}
\providecommand{\Mfield}{\mathfrak{M}}
\providecommand{\tq}{\tau}
\providecommand{\ds}{\displaystyle}

\begin{document}

\section*{\S\ 4. Рациональная база полей $\Kfield_{n\rho}$.}

В дальнейших исследованиях, сгруппированных вокруг понятий базы, мы всюду пользуемся «принципом переноса» из \S\ 3. Сначала мы доказываем существование базы для полей, но определения даем в общем виде:

\emph{Определение III: Конечное число функций из $\Ssys_{n\rho}$ называется рациональной базой, если каждую функцию из $\Ssys_{n\rho}$ можно представить как рациональную комбинацию этого конечного числа функций с коэффициентами из $\Omega$.}

Рациональная база должна содержать $\rho$ алгебраически независимых функций, поскольку алгебраический ранг системы, полученной из рациональной базы, равен алгебраическому рангу самой рациональной базы.

Сначала покажем, что доказательство существования рациональной базы допускает принцип переноса. А именно, Определение III можно сформулировать так:

Функции из $\Ssys_{n\rho}$
\[
  g_1(x),\ g_2(x),\ldots, g_k(x)
\]
тогда и только тогда образуют рациональную базу, когда выполнены бесконечно многие соотношения:
\begin{equation}
  f(x)=\gamma\bigl(g_1(x),\ldots,g_k(x)\bigr),
\tag{1}
\end{equation}
где $f(x)$ пробегает все функции из $\Ssys_{n\rho}$, а $\gamma$ обозначает рациональную функцию от $k$ аргументов, меняющуюся вместе с $f$.

Следовательно, если в системе отображения $\Ssys^*_{\rho\rho}$ рациональная база задана функциями
\[
  g_1^*(x),\ g_2^*(x),\ldots, g_k^*(x),
\]
что влечет соотношения
\[
  f^*(x)=\gamma\bigl(g_1^*(x),\ldots,g_k^*(x)\bigr),
\]
то по Теореме I для $\Ssys_{n\rho}$ выполнены соотношения (1). Иными словами, \emph{принцип переноса для рациональной базы} гласит:

\emph{Из существования рациональной базы для системы отображения $\Ssys^*_{\rho\rho}$ следует существование рациональной базы для исходной системы $\Ssys_{n\rho}$.}\footnote{Соответственно формулируется принцип переноса для каждого вопроса о базе, рассматриваемого здесь.}

Чтобы доказать существование рациональной базы всех полей $\Kfield_{n\rho}$, мы, таким образом, можем ограничиться полями $\Kfield^*_{\rho\rho}$; здесь к цели ведет прямое обобщение рассуждения, примененного Э. Штайницем.\footnote{E. Steinitz, Algebraische Theorie der Körper, \S\ 24, Crelles Journ. 137 (1909).}

Пусть
\begin{equation}
  g_1^*(x_1,\ldots,x_\rho),\ldots,
  g_\rho^*(x_1,\ldots,x_\rho)
\tag{2}
\end{equation}
есть система из $\rho$ алгебраически независимых функций из $\Kfield^*_{\rho\rho}$. Элиминация $x_1,\ldots,x_\rho$ дает $\rho$ соотношений:
\begin{equation}
\begin{aligned}
  F_1\bigl(g_1^*(x),\ldots,g_\rho^*(x),x_1\bigr)&=0,\ldots,\\
  F_\rho\bigl(g_1^*(x),\ldots,g_\rho^*(x),x_\rho\bigr)&=0,
\end{aligned}
\tag{3}
\end{equation}
причем в каждом соотношении $F_i=0$ неопределенная $x_i$ действительно входит; иначе, вопреки предположению, между функциями (2) существовала бы алгебраическая зависимость. Соотношения (3) означают, что
\[
  x_1,x_2,\ldots,x_\rho
\]
являются алгебраическими элементами относительно
\[
  \Omega\bigl(g_1^*(x),\ldots,g_\rho^*(x)\bigr).
\]
Поэтому поле, возникающее присоединением $x_1,\ldots,x_\rho$,
\[
  \Omega\bigl(g_1^*,\ldots,g_\rho^*,x_1,
  \ldots,x_\rho\bigr)=\Omega(x_1,\ldots,x_\rho),
\]
есть конечное алгебраическое расширение над $\Omega(g_1^*,\ldots,g_\rho^*)$. По известной теореме\footnote{См., например, Weber, Lehrbuch d. Algebra 1, \S\ 144 (1-е изд.) или \S\ 55 (малое изд.).} $\Kfield^*_{\rho\rho}$ как промежуточное поле между $\Omega(g_1^*,\ldots,g_\rho^*)$ и $\Omega(x_1,
\ldots,x_\rho)$ само является алгебраическим расширением над $\Omega(g_1^*,\ldots,g_\rho^*)$, тогда как $\Omega(x_1,
\ldots,x_\rho)$ является алгебраическим расширением над $\Kfield^*_{\rho\rho}$. Значит, $\Kfield^*_{\rho\rho}$ получается из $\Omega(g_1^*,\ldots,g_\rho^*)$ присоединением конечного числа элементов из $\Kfield^*_{\rho\rho}$, или также присоединением \emph{одной} подходяще выбранной функции; то есть
\[
  \Kfield^*_{\rho\rho}
  =\Omega(g_1^*,\ldots,g_\rho^*,g_{\rho+1}^*,\ldots,g_k^*)
  =\Omega(g_1^*,\ldots,g_\rho^*,g_0^*).
\]

По принципу переноса мы тем самым получаем

\emph{Теорема II: Каждое поле $\Kfield_{n\rho}$ обладает рациональной базой алгебраического ранга $\rho$; в частности, рациональную базу можно выбрать так, чтобы она содержала не более $\rho+1$ функций.}

Пусть дана произвольная система из $\rho$ алгебраически независимых функций из $\Kfield_{n\rho}$:
$h_1(x),\ldots,h_\rho(x)$. По приведенному выше методу ее можно расширить до рациональной базы
\[
  h_1(x),\ldots,h_\rho(x);
  h_{\rho+1}(x),\ldots,h_\sigma(x).
\]
По определяющему свойству существуют соотношения:
\[
\begin{aligned}
  h_i(x)&=\chi_i\bigl(g_1(x),\ldots,g_k(x)\bigr)
    &&(i=1,2,\ldots,\sigma),\\
  g_j(x)&=\psi_j\bigl(h_1(x),\ldots,h_\sigma(x)\bigr)
    &&(j=1,2,\ldots,k),
\end{aligned}
\]
где $\chi_i$, $\psi_j$ обозначают рациональные функции своих аргументов. То есть:

Из любой рациональной базы всякая другая рациональная база получается бирациональным преобразованием. Кроме того:

если
\[
  h_1(x),\ldots,h_\rho(x)
\]
есть произвольная система из $\rho$ алгебраически независимых функций из $\Kfield_{n\rho}$, то $\Kfield_{n\rho}$ является конечным алгебраическим расширением над
\[
  \Omega\bigl(h_1(x),\ldots,h_\rho(x)\bigr).
\]

\end{document}
